%=====================================================================
\section{Faktorisierung (FFF) und die Matrizen $Q_\varepsilon$}\label{sec:fff}
%=====================================================================

\subsection{Regeln von L.~Li}
Für stückweise glatte, beschränkte, periodische $f,g$ und $h=fg$:
\begin{enumerate}[nosep]
  \item keine gleichzeitigen Sprünge: $[h]=\tp f[g]$ (Laurent);
  \item nur komplementäre Sprünge ($h$ stetig): $[h]=\tp{1/f}^{-1}[g]$ (invers);
  \item gleichzeitige, nicht komplementäre Sprünge: keine der beiden Regeln.
\end{enumerate}
Die FFF zerlegt $\mathbf E$ in stetige Größen, sodass nur Fall 1 und 2 auftreten.

\subsection{Normalen- und Tangentenfelder}
$\mathbf N(x,z)$ ist ein reelles Einheitsvektorfeld, das auf jeder Materialgrenze der Schicht mit der
Flächennormalen übereinstimmt und \emph{an den Unstetigkeitsstellen von $\varepsilon$ stetig} ist
(anderswo darf es beliebig sein). Es hängt nicht von $y$ ab, wenn die Schicht durch ein Profil
$y=g(x,z)$ beschrieben wird (Normale in $(x,y,z)$ parallel zur Normale im Profilpunkt darunter).
\begin{align}
  \text{Profil }y=g(x)\ \text{(1D)}:\quad &\mathbf N=\frac{(-g',\,1,\,0)}{\sqrt{1+g'^2}},\label{eq:N1d}\\
  \text{Profil }y=g(x,z)\ \text{(2D)}:\quad &\mathbf N=\frac{(-\partial_xg,\,1,\,-\partial_zg)}{\sqrt{1+(\partial_xg)^2+(\partial_zg)^2}},\label{eq:N2d}\\
  \text{lamellar 1D}:\quad &N_x^2=1,\ N_y=N_z=0\ \text{(nur $\mathbf N\mathbf N^T$ wird gebraucht)},\label{eq:Nlam}\\
  \text{$y$-invarianter Querschnitt (2D)}:\quad &\mathbf N=(N_x,0,N_z)\ \text{aus Algorithmus \ref{alg:nfield}}.\notag
\end{align}
\begin{tcolorbox}[title=Algorithmus NF: Normalenfeld für $y$-invariante 2D-Querschnitte (v1.5)\label{alg:nfield}]
\begin{enumerate}[nosep]
  \item Raster der Einheitszelle mit $G_x\times G_z$ Punkten $x_i=d_xi/G_x$, $z_k=d_zk/G_z$,
        $G_{x,z}=\max(G,8(2N_{x,z}+1))$, Standard $G=512$.
  \item \textbf{Materialgrenzen} exakt aus den Polygonen: jede Kante wird in 64 Stücke geteilt; ein Stück ist
        Grenze, wenn die Materialien beidseits (periodische Punktsuche, Abstand $10^{-9}d$) verschieden sind.
        Damit entfallen Kanten auf dem Zellrand, die an ein periodisches Bild desselben Materials stoßen,
        innere Kanten aneinanderstoßender Einschlüsse gleichen Materials und Einschlüsse aus Hintergrundmaterial.
        Für jeden Rasterpunkt $\mathbf r$: nächster Grenzpunkt $\mathbf p$ (Grenzstücke mit Bildern $\pm d_x,\pm d_z$),
        Richtung $\mathbf n=(\mathbf r-\mathbf p)/|\mathbf r-\mathbf p|$ (auf der Grenze: Kantennormale).
        Gleich weit entfernte Kandidaten (Mittelachse, Toleranz $10^{-12}d$) werden gemittelt (Symmetrie).
  \item \textbf{Glättung} mit periodischem Gaußfilter gleicher \emph{physikalischer} Breite in $x$ und $z$
        ($\sigma$ Rasterpunkte entlang der kürzeren Periode, Standard 3).
  \item \textbf{Tensorvariante (Standard):} geglättet wird das Projektorfeld $\mathbf n\mathbf n^T$; das Ergebnis
        ist direkt $\tp{\cdot}$-Quelle für $N_iN_j$ (symmetrisch, Spur 1, stetig, an geraden Grenzen $=\mathbf n\mathbf n^T$).
        \textbf{Vektorvariante} (Fassung v1.4): geglättet wird $\nabla\phi$ der vorzeichenbehafteten Abstandsfunktion
        ($\phi>0$ im Hintergrund), $\mathbf N=\mathbf g/|\mathbf g|$ wo $|\mathbf g|>\tau\max|\mathbf g|$ ($\tau=10^{-3}$),
        sonst $N_xN_x=N_zN_z=\tfrac12$, $N_xN_z=0$.
  \item Fourierkoeffizienten von $N_x^2$, $N_xN_z$, $N_z^2$ per DFT des Rasters (glatte Funktionen).
\end{enumerate}
Literatur: Schuster et al.\ (JOSA A 24, 2880, 2007); Götz et al.\ (Opt.\ Express 16, 17295, 2008).
\end{tcolorbox}
\textbf{Begründung der Tensorvariante (Befund M7).} Für die FFF-Herleitung \eqref{eq:Qiso} genügt
\emph{irgendeine} stetige Zerlegung $\mathbf E=(\mathbb 1-P)\mathbf E+P\mathbf E$, deren Matrixfeld $P$ an den
Unstetigkeitsstellen von $\varepsilon$ gleich $\mathbf n\mathbf n^T$ ist; $P$ muss weder Projektor noch aus einem
Vektorfeld gebildet sein. An Mittelachsen zwischen gegenüberliegenden Wänden (Lamellenmitte, Steg zwischen zwei
Löchern) kehrt $\nabla\phi$ das Vorzeichen um; die Vektorglättung läuft dort durch null und fällt auf $\tfrac12$
zurück. Liegt die Mittelachse auf Rasterpunkten, weicht $N_x^2$ einer $z$-invarianten Lamelle in den Koeffizienten um
$\approx2\cdot10^{-3}$ von 1 ab (Unit-Test \texttt{VectorMethodFallbackOnMedialAxis}); V9 wäre nicht exakt. Die
Tensorvariante ist vorzeichenfrei, singularitätsfrei und reproduziert Lamellen exakt ($10^{-14}$). Sie ist
auch für die Projektorform \eqref{eq:QaniB} (Variante B, M8) direkt verwendbar. An konvexen Ecken liefert sie
gemischte (Rang-2-)Tensoren statt gedrehter Normalen; beide Varianten weichen dort von $\mathbf n\mathbf n^T$ ab.
Ein stetiges \emph{vorzeichenbehaftetes} Feld existiert für Löcher nicht; \texttt{FFF\_rig} ist für
\texttt{Polygon2D} daher nicht zulässig.
Tangentenvektoren für Variante~A (unten) mit $\mathbf T_2\times\mathbf T_1=\mathbf N$:
\begin{align}
  \text{Profilschichten }(N_y>0):\quad&\mathbf T_1=\frac{(N_y,-N_x,0)}{\sqrt{N_x^2+N_y^2}},\quad
  \mathbf T_2=\mathbf T_1\times\mathbf N=\frac{(-N_xN_z,\,-N_yN_z,\,N_x^2+N_y^2)}{\sqrt{N_x^2+N_y^2}},\label{eq:T_profile}\\
  \text{$y$-invariante Schichten }(N_y=0):\quad&\mathbf T_1=\hat{\mathbf y},\quad \mathbf T_2=\hat{\mathbf y}\times\mathbf N=(N_z,0,-N_x).\label{eq:T_yinv}
\end{align}
(Prüfung: $(N_z,0,-N_x)\times(0,1,0)=(N_x,0,N_z)=\mathbf N$. Die Buchwahl \eqref{eq:T_profile} ist für
$N_x=N_y=0$ singulär und daher für $y$-invariante Schichten ungeeignet.)

\subsection{Isotrope Materialien}
Mit $\Delta=\tp{\varepsilon}-\Li$ (Buch (IV.32), (IX.11)):
\begin{equation}
  Q_{\varepsilon}=\tp{\varepsilon}\otimes\mathbb 1_3-\Delta\,\tp{\mathbf N\mathbf N^T},\qquad
  Q_{\varepsilon,ij}=\delta_{ij}\tp{\varepsilon}-\Delta\tp{N_iN_j}.
  \label{eq:Qiso}
\end{equation}
Herleitung: $\mathbf D=\varepsilon\mathbf E_T+\varepsilon\mathbf E_N$; $\varepsilon\mathbf E_T$: Laurent
($\mathbf E_T$ stetig), $\varepsilon\mathbf E_N=\mathbf D_N$ stetig: invers; mit $\mathbf E_N=\mathbf N\mathbf N^T\mathbf E$
und $\mathbf E_T=\mathbf E-\mathbf E_N$ folgt \eqref{eq:Qiso} (Laurent für $\mathbf N\mathbf N^T\mathbf E$, da $\mathbf N$ stetig).
Grenzfälle (Unit-Tests): $N\to\pm\hat x$ überall $\Rightarrow$ Li-Formulierung für Lamellen
($Q_{xx}=\Li$, $Q_{yy}=Q_{zz}=\tp\varepsilon$); $\mathbf N=\hat y$ überall $\Rightarrow$ klassische Laurent-Formulierung
($Q_{yy}=\Li$, sonst $\tp\varepsilon$, vgl.\ Buch IV.2.3: $a/d\to0$).

\subsection{Varianten der isotropen FFF (Befund aus M3)}\label{sec:fffvariants}
Die Buchformel \eqref{eq:Qiso} entsteht aus einer exakteren Form durch zwei Näherungen im
abgeschnittenen Raum. Aus $D_i=\varepsilon E_{T,i}+N_iD_N$ mit $E_{T,i}=E_i-N_iE_N$ (stetig), $D_N$ (stetig),
$E_N=\sum_jN_jE_j$ folgt mit den Li-Regeln (Laurent für $\varepsilon\cdot$stetig und $N\cdot$beliebig,
invers für $D_N=\varepsilon E_N$):
\begin{equation}
  Q^{(\mathrm{rig})}_{ij}=\delta_{ij}\tp{\varepsilon}-\tp{\varepsilon}\tp{N_i}\tp{N_j}+\tp{N_i}\,\Li\,\tp{N_j}.
  \label{eq:Qrig}
\end{equation}
Das Buch ersetzt $\tp{N_i}\tp{N_j}\to\tp{N_iN_j}$ und vertauscht $\tp{N_i}$ mit $\Li$; das ergibt
\eqref{eq:Qiso}. Eine dritte Variante symmetrisiert das Produkt:
\begin{equation}
  Q^{(\mathrm{sym})}_{ij}=\delta_{ij}\tp{\varepsilon}-\tfrac12\big(\Delta\tp{N_iN_j}+\tp{N_iN_j}\Delta\big).
  \label{eq:Qsym}
\end{equation}
\textbf{Numerischer Befund} (RK, TM, $\lambda=0{,}6328\,\mu$m; Referenzen aus der C-Methode):
\begin{center}
\small
\begin{tabular}{lcc}
\toprule
& Gold-Sinus $d=1$, $a=0{,}25\,\mu$m, $20^\circ$, $e_{-1}$ ($N=8/16/24$) & Al-Sinus V4, $e_{-1}$ ($N=20$) \\
\midrule
Referenz & $0{,}26614$ & $0{,}8696$ \\
\verb|FFF_iso| \eqref{eq:Qiso} & falsch ($0{,}327$ / $0{,}217$; $\sum e>1$) & $0{,}86951$ \\
\verb|FFF_sym| \eqref{eq:Qsym} & $0{,}2681$ / $0{,}2672$ / $0{,}2665$ & $0{,}86565$ \\
\verb|FFF_rig| \eqref{eq:Qrig} & $0{,}26609$ / $0{,}26621$ / -- & $0{,}87075$ \\
\bottomrule
\end{tabular}
\end{center}
\begin{itemize}[nosep]
  \item Die Buchvariante versagt bei schwach absorbierenden Metallen ($\operatorname{Re}\varepsilon<0$, kleines
        $\operatorname{Im}\varepsilon$). Vermuteter Mechanismus: $1/\varepsilon$ läuft fast durch null, $\Li$ wird groß,
        und damit auch der Fehler der Vertauschung von $\tp{N_i}$ und $\Li$. Der Löser gibt dann eine Warnung aus.
  \item In einer $y$-invarianten Scheibe aus isotropem Material müssen die Eigenwerte als $\pm\lambda$-Paare
        auftreten. Laurent und \verb|FFF_sym| erhalten diese Paarung exakt, \verb|FFF_iso| und \verb|FFF_rig|
        nicht. Da \verb|FFF_rig| trotzdem am genauesten ist, ist die Paarung \emph{kein} Gütekriterium.
  \item \verb|FFF_rig| braucht das \emph{vorzeichenbehaftete} Normalenfeld; springt es (Ecken von Dreiecks-
        und Echelette-Profilen), versagt die Variante ($\sum e>1$). Dort ist \verb|FFF_sym| zu verwenden.
  \item Keine der FFF-Varianten erhält bei verlustfreien Medien die Energie exakt (Abweichung
        $10^{-5}$--$5\cdot10^{-4}$), Laurent schon. Tests prüfen die Energie bei Profilen daher mit Toleranz.
\end{itemize}
\textbf{2D (M7):} Für \texttt{Polygon2D} (Querschnitte mit Löchern/Säulen) ist \verb|FFF_sym| Standard;
\verb|FFF_iso| liefert beim Ag-Lochgitter (V10, $\varepsilon=-79{,}9+1{,}79i$) bei $N=3$ und $N=6$ $R>1$
(dieselbe Ursache wie bei Gold in 1D), \verb|FFF_rig| ist mangels stetigem Vorzeichenfeld nicht anwendbar.
\texttt{Profile2D} verwendet wie in 1D \verb|FFF_rig| (analytische Formen) bzw.\ \verb|FFF_sym| (Funktionen);
$\mathbf N$ aus \eqref{eq:N2d} ist dort stetig und vorzeichenbehaftet.

\textbf{Standard (\texttt{Factorization::Auto})} für isotrope Profilschichten: \verb|FFF_rig| bei glatten
Profilen (Sinus), \verb|FFF_sym| bei Profilen mit Ecken (stückweise linear) und bei benutzerdefinierten
Funktionen. \verb|FFF_iso| bleibt als Buchvariante wählbar. Beide neuen Varianten sind Eigenherleitungen
dieses Projekts; ob sie so in der Literatur stehen, ist noch zu prüfen. Offener Punkt: Für Gold-Dreiecke
gibt es noch keine belastbare Referenz (V14).

\subsection{Anisotrope Materialien, Variante A (Buch Kap.~V)}
Stetiger Vektor $\mathbf F_\varepsilon=(E_{T_1},D_N,E_{T_2})=A_\varepsilon\mathbf E$ mit
\begin{equation}
  A_\varepsilon=\begin{pmatrix}\mathbf T_1^T\\ \mathbf v^T\\ \mathbf T_2^T\end{pmatrix},\quad
  \mathbf v^T=\mathbf N^T\overline{\overline\varepsilon}\ \ \big(v_j=\textstyle\sum_iN_i\varepsilon_{ij}\big),\qquad
  \xi=\mathbf N^T\overline{\overline\varepsilon}\mathbf N,
  \label{eq:Aeps}
\end{equation}
\begin{equation}
  C_\varepsilon=A_\varepsilon^{-1}=\frac{1}{\xi}\Big(\ \mathbf v\times\mathbf T_2\ \Big|\ \mathbf N\ \Big|\ -\mathbf v\times\mathbf T_1\ \Big)
  \quad\text{(Spalten)},\qquad \det A_\varepsilon=\xi .
  \label{eq:Ceps}
\end{equation}
(Buchtext S.\,102 nennt die Spaltenreihenfolge falsch; (V.11) selbst ist richtig.)
$\mathbf D=\overline{\overline\varepsilon}C_\varepsilon\mathbf F_\varepsilon$: Laurent; $\mathbf F_\varepsilon=A_\varepsilon\mathbf E$ stetig: invers.
\begin{equation}
  Q_\varepsilon=\tp{\overline{\overline\varepsilon}C_\varepsilon}\,\tp{C_\varepsilon}^{-1}
  \label{eq:QaniA}
\end{equation}
($3\times3$ Blöcke; $\tp{X}$ für eine matrixwertige Funktion $X$ ist die Blockmatrix der Toeplitz-Matrizen
der Einträge). Speziell 1D-Profil ($\mathbf T_2=\hat z$, $\mathbf T_1=(N_y,-N_x,0)$):
\begin{equation}
  C_\varepsilon=\frac1\xi\begin{pmatrix}v_y&N_x&-v_zN_x\\-v_x&N_y&-v_zN_y\\0&0&\xi\end{pmatrix},\qquad
  \xi=v_xN_x+v_yN_y
  \label{eq:C1d}
\end{equation}
(Buch (VIII.2) hat in $C_{23}$ ein falsches Vorzeichen). \textbf{Voraussetzung:} $\xi\ne0$ überall
(bei indefiniten, \enquote{hyperbolischen} Tensoren nicht garantiert; dann Warnung und Variante B
oder Laurent). \textbf{Achtung:} $C_\varepsilon$ hängt über $\overline{\overline\varepsilon}(x,y,z)$ von $y$ ab, außer bei
$y$-invarianten Schichten -- $\tp{C_\varepsilon}$ muss für jede Höhe neu berechnet werden (Buch S.\,108
behauptet Gegenteiliges).

\paragraph{Lamellen mit Variante A.} Für 1D-Lamellen genügt das \emph{konstante} Feld $\mathbf N=\hat x$:
An jeder Wand ist es bis auf das Vorzeichen die Normale, und das Vorzeichen ist überall gleich, also ist
$\mathbf F_\varepsilon$ stetig. Mit isotropem Tensor reproduziert Variante A dann exakt die Li-Lamellenformel
(V8, $10^{-10}$). Standard für anisotrope Lamellen und Profile ist \verb|FFF_A|.
\paragraph{Fourierkoeffizienten.} $C_\varepsilon(x)$ und $\overline{\overline\varepsilon}C_\varepsilon(x)$ sind stückweise glatt
(Tensor stückweise konstant, $\mathbf N$ glatt). Koeffizienten per zusammengesetzter 8-Punkt-Gauß-Quadratur je
Intervall mit $|pK|h\le1$ (\verb|piecewiseSmoothCoefficients|); für konstante Funktionen exakt bis auf Rundung.

\subsection{Anisotrope Materialien, Variante B (vorzeichenunabhängig)}
Variante A verlangt ein stetiges $\mathbf N$ \emph{mit Vorzeichen}. Bei 2D-Querschnitten (Löcher, Säulen)
gibt es kein solches Feld; dort würden $E_{T_1}$, $D_N$ im Inneren eines homogenen Gebiets springen. Variante B
verwendet nur die Projektoren $P_N=\mathbf N\mathbf N^T$, $P_T=\mathbb 1-P_N$ und die stetigen
\emph{Vektoren} $\mathbf e_t=P_T\mathbf E$, $\mathbf d_n=P_N\mathbf D$:
aus $\mathbf N^T\mathbf D=\mathbf N^T\overline{\overline\varepsilon}\mathbf e_t+\xi\,\mathbf N^T\mathbf E$ folgt
\begin{equation}
  \mathbf E=\Big(\mathbb 1-\tfrac1\xi P_N\overline{\overline\varepsilon}\Big)\mathbf e_t+\tfrac1\xi\mathbf d_n,\qquad
  \mathbf D=\overline{\overline\varepsilon}\Big(\mathbb 1-\tfrac1\xi P_N\overline{\overline\varepsilon}\Big)\mathbf e_t+\tfrac1\xi\overline{\overline\varepsilon}\,\mathbf d_n .
  \label{eq:projector}
\end{equation}
Alle Koeffizientenfunktionen sind invariant unter $\mathbf N\to-\mathbf N$. Laurent auf beide
Zeilen, $[\mathbf e_t]=\tp{P_T}[\mathbf E]$ (Laurent), Elimination von $[\mathbf d_n]$:
\begin{equation}
  Q_\varepsilon^{(B)}=\tp{G_1}\tp{P_T}+\tp{G_2}\,\tp{\xi^{-1}}^{-1}\big(\mathbb 1-\tp{H_1}\tp{P_T}\big),
  \label{eq:QaniB}
\end{equation}
\[
  G_1=\overline{\overline\varepsilon}\big(\mathbb 1-\xi^{-1}P_N\overline{\overline\varepsilon}\big),\quad
  G_2=\xi^{-1}\overline{\overline\varepsilon},\quad H_1=\mathbb 1-\xi^{-1}P_N\overline{\overline\varepsilon},
\]
mit $\tp{\xi^{-1}}$ = blockdiagonal ($3\times$ dieselbe skalare Toeplitz-Matrix). Für isotropes $\varepsilon$
reduziert sich \eqref{eq:QaniB} bis auf Abschneideeffekte ($\tp{P_T}\tp{P_T}\approx\tp{P_T}$) auf
\eqref{eq:Qiso}. \textbf{Test:} Variante A und B müssen auf Profilschichten mit $N$ gegen dieselben
Effizienzen konvergieren; Variante B mit isotropem Tensor gegen \eqref{eq:Qiso}.
\emph{Hinweis:} Variante B ist eine Eigenherleitung dieses Plans (nicht im Buch). M5-Befund: am Kobalt-Gitter
(V7) stimmt sie mit Variante A auf 1--3\,\% (TM-Anteile) bzw.\ 2--15\,\% (sehr kleine TE-Anteile) überein; ein
Literaturvergleich (Li, J.\ Mod.\ Opt.\ 45, 1313 (1998)) steht noch aus. Sie ist für 2D-Querschnitte (M8) gedacht.

\paragraph{Variante B für \texttt{Polygon2D} (M8, v1.6).} $P_N$ muss an den Grenzflächen ein Projektor
$\mathbf N\mathbf N^T$ sein und $\xi=\mathbf N^T\overline{\overline\varepsilon}\mathbf N$ bleibt nur für Rang-1-Felder
sinnvoll; das Tensorfeld von Algorithmus NF (Rang 2 an Ecken) ist hier ungeeignet (für Metalle könnte
$\mathbb 1-P+P\overline{\overline\varepsilon}$ singulär werden). Verwendet wird das \emph{Rang-1-Feld}
$\mathbf N=\mathbf g/|\mathbf g|$ der geglätteten vorzeichenbehafteten Abstandsfunktion (Vektorvariante von
Algorithmus NF; Vorzeichenwechsel an Mittelachsen sind für B bedeutungslos), wo $|\mathbf g|\le\tau\max|\mathbf g|$
die Hauptachse des geglätteten Tensors (z.\,B.\ Lamellenmitte: exakt $\hat x$). Die Koeffizientenfunktionen
$F\in\{G_1,G_2,H_1,P_T,\xi^{-1}\}$ springen mit dem Material:
\[
  F=F(\overline{\overline\varepsilon}_{\text{bg}},\mathbf N)+\textstyle\sum_k\chi_k\big[F(\overline{\overline\varepsilon}_k,\mathbf N)-F(\overline{\overline\varepsilon}_{\text{bg}},\mathbf N)\big].
\]
Glatte Anteile per DFT des NF-Rasters (Rastergröße Zweierpotenz, FFT), Produkte mit $\chi_k$ per Faltung
$\sum_{p'q'}\hat D_{p'q'}\hat\chi_{p-p',q-q'}$ mit den exakten Polygonkoeffizienten und $P_f=2P$ Ordnungen des
glatten Faktors (Laurent-Regel für glatt $\times$ unstetig). Für konstantes $\mathbf N$ exakt: ein
$z$-invarianter Streifen liefert mit isotropem Tensor die Li-Formel ($10^{-10}$, V8-2D). Voraussetzung:
Einschlüsse liegen im Hintergrund (das Vorzeichen von $\phi$ ist sonst nicht definiert); sich berührende
Einschlüsse verschiedener Materialien nur mit Laurent. Standard: anisotrope \texttt{Polygon2D} \verb|FFF_B|,
anisotrope \texttt{Profile2D} \verb|FFF_A| ($\mathbf N$ aus \eqref{eq:N2d} stetig und vorzeichenbehaftet,
$\mathbf T_1=(N_y,-N_x,0)/\sqrt{N_x^2+N_y^2}$, $\mathbf T_2=\mathbf T_1\times\mathbf N$).

\subsection{Weitere Faktorisierungsoptionen (für Vergleiche)}
\verb|Laurent| ($Q_{ij}=\tp{\varepsilon_{ij}}$), \verb|LiLamellar| (nur 1D-Lamellen, Buch (IV.5)/(IV.6)),
\verb|FFF_iso| \eqref{eq:Qiso}, \verb|FFF_sym| \eqref{eq:Qsym}, \verb|FFF_rig| \eqref{eq:Qrig},
\verb|FFF_A| \eqref{eq:QaniA}, \verb|FFF_B| \eqref{eq:QaniB}.
Die Option wird je Schicht gewählt; Standard: isotrope Lamellen \verb|LiLamellar|, isotrope Profile (1D und 2D) siehe
Abschnitt~\ref{sec:fffvariants}, isotrope \texttt{Polygon2D} \verb|FFF_sym|, anisotrop Profil (1D und 2D) \verb|FFF_A|,
anisotrop $y$-invariant (Lamellen: \verb|FFF_A| mit $\mathbf N=\hat x$; \texttt{Polygon2D}) \verb|FFF_B|.

%=====================================================================
\section{Homogene Medien: Modenbasen}\label{sec:homog}
%=====================================================================
In homogenen Medien (Substrat, Superstrat, homogene Zwischenschichten, Referenzmedium der
dünnen Schichten) ist $M$ \emph{ordnungsweise} entkoppelt: $M$ zerfällt in $n_f$ unabhängige
$4\times4$-Matrizen $M^{(r)}$ (eine je Ordnung $r\leftrightarrow(n,m)$).

\subsection{Isotrop: analytische Basis der tangentialen E-Komponenten}
Für eine ebene Welle $e^{i(\alpha x+k_yy+\gamma z)}$ ($k_y=+\beta$ auf, $-\beta$ ab) gilt in normierten
Einheiten $\mathbf H=\mathbf k\times\mathbf E$ ($\mu=1$), $\mathbf k\cdot\mathbf E=0$. Mit
$E_y=-(\alpha E_x+\gamma E_z)/k_y$ und $\alpha^2+k_y^2+\gamma^2=\varepsilon$:
\begin{equation}
  \begin{pmatrix}H_x\\H_z\end{pmatrix}=\frac{1}{k_y}K\begin{pmatrix}E_x\\E_z\end{pmatrix},\qquad
  K=\begin{pmatrix}\alpha\gamma&\varepsilon-\alpha^2\\-(\varepsilon-\gamma^2)&-\alpha\gamma\end{pmatrix},\qquad
  \det K=\varepsilon\beta^2 .
  \label{eq:Kmat}
\end{equation}
Amplitudenvektor $V=(\mathbf a;\mathbf b)$ mit $\mathbf a=(a_x,a_z)$ (ab), $\mathbf b=(b_x,b_z)$ (auf) = tangentiale
E-Komponenten der Teilwellen am Bezugspunkt. Mit $\mathcal V=K\beta^{-1}$ (ordnungsweise $2\times2$):
\begin{equation}
  F=\Psi V,\qquad \Psi=\begin{pmatrix}\mathbb 1&\mathbb 1\\-\mathcal V&\mathcal V\end{pmatrix},\qquad
  \Psi^{-1}=\frac12\begin{pmatrix}\mathbb 1&-\mathcal V^{-1}\\\mathbb 1&\mathcal V^{-1}\end{pmatrix},\qquad
  \mathcal V^{-1}=\frac{1}{\varepsilon\beta}\begin{pmatrix}-\alpha\gamma&-(\varepsilon-\alpha^2)\\\varepsilon-\gamma^2&\alpha\gamma\end{pmatrix}.
  \label{eq:Psi_iso}
\end{equation}
Die Spalten sind Eigenvektoren von $M^{(r)}$ zu $\lambda=-i\beta$ (ab) und $\lambda=+i\beta$ (auf).
(Diese Basis ist der in Buch Anh.~IV vorzuziehen: dort wird durch $k^2-\gamma_m^2$ dividiert, was im 2D-Fall
verschwinden kann; hier tritt nur $\beta$ im Nenner auf.) Kontrollen: 1D-TE ($\gamma=0$, $E_x=0$):
$H_x=k_yE_z$; 1D-TM: $E_x=-k_yH_z/\varepsilon$.

\subsection{Anisotrop: $4\times4$-Eigenzerlegung je Ordnung (Berreman)}
$M^{(r)}$ aus \eqref{eq:Mgeneral} mit $Q_{\varepsilon,ij}=\varepsilon_{ij}$ (Skalare) und $\alpha_n,\gamma_m$.
Eigenwerte $\lambda_k$, Eigenvektoren $w_k$ ($k=1..4$). Klassifikation (Feld $\propto e^{\lambda y}$):
\begin{equation}
  \text{auf (b)}:\ \operatorname{Re}\lambda<-\tau_\lambda\ \ \text{oder}\ \ \big(|\operatorname{Re}\lambda|\le\tau_\lambda\ \text{und}\ S_y(w)>0\big),
  \qquad\text{sonst ab (a)},
  \label{eq:classify}
\end{equation}
mit dem zeitgemittelten Poynting-Fluss der Mode
\begin{equation}
  S_y(w)=\tfrac12\operatorname{Re}\big(E_zH_x^*-E_xH_z^*\big)
  \label{eq:Sy}
\end{equation}
($E_x,E_z,H_x,H_z$ = Komponenten von $w$). Es müssen genau zwei auf- und zwei abwärts
klassifizierte Moden entstehen, sonst Fehler. Normierung: $\Psi$-Spalten so skalieren, dass
$\max(|E_x|,|E_z|)=1$. Für isotrope Medien muss dieses Verfahren \eqref{eq:Psi_iso} (bis auf
Skalierung/Mischung entarteter Moden) reproduzieren (Unit-Test).

\subsection{Leistungsfluss und Effizienzen}
In einem homogenen Medium sind die Ordnungen bezüglich des Flusses orthogonal; die von der Ordnung
$r$ in $\pm y$ transportierte Leistung (pro Periode, normiert) ist \eqref{eq:Sy} mit den Fourierkoeffizienten
dieser Ordnung. Einfallende Leistung $S_{\text{inc}}=|S_y(\text{einf.\ Welle})|$. Effizienzen:
\begin{equation}
  e^{(r)}_{nm}=\frac{S_y^{\uparrow}(r)}{S_{\text{inc}}},\qquad e^{(t)}_{nm}=\frac{|S_y^{\downarrow}(r)|}{S_{\text{inc}}}
  \quad(\text{nur propagierende Ordnungen in verlustfreien Medien}).
  \label{eq:eff}
\end{equation}
Geschlossene Form für isotrope verlustfreie Medien (Buch (IX.20), korrigiert; normiert,
$k^2\to\varepsilon$, $|\mathbf A|=1$):
\begin{equation}
  e_{nm}=\frac{(\varepsilon-\gamma_m^2)|B_x|^2+(\varepsilon-\alpha_n^2)|B_z|^2+2\alpha_n\gamma_m\operatorname{Re}(B_x\overline{B_z})}
  {\beta_{00}^{(M)}\,\beta_{nm}},
  \label{eq:eff_closed}
\end{equation}
wobei $\varepsilon$ und $\beta_{nm}$ zu dem Medium gehören, in dem die Ordnung läuft (für die
Transmission also $\varepsilon_1$, $\beta^{(1)}_{nm}$; das Buch vergisst $k_2\to k_1$), und
$\beta^{(M)}_{00}=\beta_0$. Bei absorbierendem Substrat nur Reflexion und Absorption $1-\sum e^{(r)}$.
Für die Ausgabe zusätzlich s/p-Zerlegung je Ordnung: mit $\kappa=(\alpha_n,\gamma_m)$,
$\hat{\mathbf s}=(\kappa_z,0,-\kappa_x)/|\kappa|$ (für $\kappa=0$: $\hat{\mathbf s}=\hat z$),
$\hat{\mathbf p}=\hat{\mathbf k}\times\hat{\mathbf s}$; Amplituden $\mathbf E\cdot\hat{\mathbf s}$, $\mathbf E\cdot\hat{\mathbf p}$
mit $E_y$ aus $\mathbf k\cdot\mathbf E=0$.

%=====================================================================
\section{S-Matrix-Formalismus}\label{sec:smatrix}
%=====================================================================

\subsection{Konventionen}
Zwischen zwei benachbarten Schichten liegt jeweils eine \emph{unendlich dünne Schicht aus dem
Referenzmedium} (Standard: Superstrat; konfigurierbar). Dort werden die Felder durch
$V=(\mathbf a;\mathbf b)$ in der Basis $\Psi_{\text{ref}}$ \eqref{eq:Psi_iso} beschrieben, mit dem Phasenbezug
\emph{lokal an dieser Grenzfläche}. Ein Block (Schicht, Schichtfolge) mit unterer Grenze \enquote{u}
und oberer Grenze \enquote{o} hat die S-Matrix
\begin{equation}
  \begin{pmatrix}\mathbf b_{\text{o}}\\\mathbf a_{\text{u}}\end{pmatrix}=
  \begin{pmatrix}S_{11}&S_{12}\\S_{21}&S_{22}\end{pmatrix}
  \begin{pmatrix}\mathbf b_{\text{u}}\\\mathbf a_{\text{o}}\end{pmatrix}
  \label{eq:Sdef}
\end{equation}
(auslaufend = $S\cdot$einlaufend; $S_{11}$: Transmission aufwärts, $S_{12}$: Reflexion von oben,
$S_{21}$: Reflexion von unten, $S_{22}$: Transmission abwärts; Buch (III.28)). Blockgröße
$n_h=2n_f$ (zwei Polarisationskomponenten je Ordnung; im 1D-in-plane-Fall $n_h=n_f$).

\subsection{Transfermatrix $\to$ S-Matrix}
Ist $T$ die Transfermatrix $(\mathbf a_{\text{o}};\mathbf b_{\text{o}})=T(\mathbf a_{\text{u}};\mathbf b_{\text{u}})$, so gilt
(Buch (III.30), (III.31), ergänzt):
\begin{equation}
  S_{11}=T_{22}-T_{21}T_{11}^{-1}T_{12},\quad S_{12}=T_{21}T_{11}^{-1},\quad
  S_{21}=-T_{11}^{-1}T_{12},\quad S_{22}=T_{11}^{-1}.
  \label{eq:T2S}
\end{equation}
Nur für dünne Scheiben verwenden (Abschnitt~\ref{sec:rk}); $T_{11}$ mit LU faktorisieren, nicht invertieren.

\subsection{Redheffer-Sternprodukt}
Für einen unteren Block $A$ und einen darüberliegenden Block $B$ (gemeinsame innere Grenzfläche mit
$\mathbf a_m,\mathbf b_m$) folgt aus $\mathbf b_m=A_{11}\mathbf b_{\text{in}}+A_{12}\mathbf a_m$,
$\mathbf a_m=B_{21}\mathbf b_m+B_{22}\mathbf a_{\text{in}}$:
\begin{equation}
\begin{aligned}
  D_1&=(\mathbb 1-A_{12}B_{21})^{-1}, & D_2&=(\mathbb 1-B_{21}A_{12})^{-1},\\
  C_{11}&=B_{11}D_1A_{11}, & C_{12}&=B_{12}+B_{11}D_1A_{12}B_{22},\\
  C_{21}&=A_{21}+A_{22}D_2B_{21}A_{11}, & C_{22}&=A_{22}D_2B_{22},
\end{aligned}
\qquad C=A\star B .
\label{eq:star}
\end{equation}
Neutrales Element: $S_{11}=S_{22}=\mathbb 1$, $S_{12}=S_{21}=0$. Das Buch-Rekursionsschema
(III.44)--(III.47) ist der Spezialfall \enquote{Stapel $\star$ Scheibe aus $T$}; es enthält in (III.45) einen Fehler
(Buch: $S_{21}$ statt $S_{22}$ im zweiten Term) -- im Code ausschließlich \eqref{eq:star} verwenden.
\textbf{Test:} Sternprodukt beliebiger Folgen = S aus dem Produkt der $T$-Matrizen (für gut konditionierte $T$).

\subsection{Grenzflächen und Gesamtstapel}
Übergang von einem homogenen Medium mit Basis $\Psi_X$ in die dünne Referenzschicht (Dicke 0):
$T=\Psi_{\text{ref}}^{-1}\Psi_X$ bzw.\ umgekehrt; S per \eqref{eq:T2S} (bei Dicke 0 unkritisch).
Diese Matrix ersetzt die Buchformeln (III.24)/(III.26) (dort sind $C$ und $C'$ vertauscht) und behebt
das Buchproblem $A'^{(1)}=A^{(1)}$ (gilt nur bei $\nu_1=\nu_M$). Gesamtstapel:
\begin{equation}
  S_{\text{tot}}=S_{\text{Substrat}\to\text{ref}}\star S_{\text{Schicht }2}\star\cdots\star S_{\text{Schicht }M-1}\star S_{\text{ref}\to\text{Superstrat}} .
  \label{eq:Stot}
\end{equation}
Mit Einfall nur von oben ($\mathbf b_{\text{Substrat,in}}=0$) und $\mathbf a_{\text{in}}=(A_x,A_z)$ in der Ordnung $(0,0)$:
\begin{equation}
  \mathbf b_{\text{Superstrat}}=S_{\text{tot},12}\,\mathbf a_{\text{in}},\qquad
  \mathbf a_{\text{Substrat}}=S_{\text{tot},22}\,\mathbf a_{\text{in}} .
  \label{eq:solution}
\end{equation}
Die Spalte der Ordnung $(0,0)$ genügt (\enquote{zentrale Spalte}); für Einfall aus mehreren Richtungen
(z.\,B.\ Reziprozitätstest) beliebige Spalten.

\subsection{Interne Amplituden (für das Nahfeld)}
Gespeichert werden für jede innere Grenzfläche $q$ die kumulierte S-Matrix des Teilstapels darunter,
$S_L^{(q)}$, und darüber, $S_U^{(q)}$. Mit den bekannten Eingängen $\mathbf b_{\text{in}}$ (unten, meist 0)
und $\mathbf a_{\text{in}}$ (oben):
\begin{equation}
  \big(\mathbb 1-S^{(q)}_{L,12}S^{(q)}_{U,21}\big)\mathbf b_q=S^{(q)}_{L,11}\mathbf b_{\text{in}}+S^{(q)}_{L,12}S^{(q)}_{U,22}\mathbf a_{\text{in}},\qquad
  \mathbf a_q=S^{(q)}_{U,21}\mathbf b_q+S^{(q)}_{U,22}\mathbf a_{\text{in}} .
  \label{eq:internal}
\end{equation}
Speicherbedarf: $2\cdot4n_h^2\cdot16$\,Byte je Grenzfläche (2D, $N=10$: $n_h=882$, ca.\ 100\,MB je
Grenzfläche). Alternative bei Speichermangel: nur $S_L$ speichern und $S_U$ in einem zweiten Durchlauf
von oben erzeugen.
